\subsection{Notation and convention}
\begin{itemize}\itemsep=0pt

\item Unless stated otherwise, categories and stacks we consider are over an algebraically closed field $k$ of characteristic zero.

\item For a character $\chi\colon G\to \bG_m$ of an algebraic group $G$, we denote by $\cO_X(\chi)$, or~sim\-ply~$\cO(\chi)$, the $G$-equivariant invertible sheaf on a $G$-scheme $X$ associated to $\chi$.

\item For an exact category $\scrE$, we denote by $\Ch(\scrE)$ (resp.\ $\Chb(\scrE)$) the category of cochain complexes (resp.\ bounded cochain complexes) in $\scrE$, and we write $\D(\scrE)$ (resp.\ $\Db(\scrE)$) for~the derived category (resp.\ the bounded derived category) of $\scrE$.
\end{itemize}

\section{Preliminaries}%\label{preliminaries}
In this section, we recall definitions and basic properties about derived factorization categories, equivariant quasi-coherent sheaves and tilting objects.

\subsection{Derived factorization categories}%\label{section 2.1}
We recall the basics of derived factorization categories mainly to fix notation. See, for example,~\cite{bdfik,H1,posi} for more details.
Throughout this section, $\scrA$ is an abelian category with small coproducts such that the small coproducts of families of short exact sequences are exact. In~what follows, we fix a triple
\begin{equation}\label{triple}
(\scrE,\Phi,w)
\end{equation}
consisting of an exact subcategory $\scrE\subseteq\scrA$ of $\scrA$, an exact autoequivalence $\Phi\colon \scrA\to \scrA$ pre\-ser\-ving~$\scrE$ and a functor morphism $w\colon \id\to\Phi$ that is compatible with $\Phi$, i.e., the equality $w({\Phi(E)})=\Phi(w(E))$ holds for every object $E\in\scrE$. The functor morphism $w\colon \id\to\Phi$ is called a {\it potential} of $\scrE$.

\begin{dfn} A {\it factorization} of $(\scrE,\Phi,w)$ is a sequence
\begin{gather*}
E=\Bigl(E_1\xrightarrow{\varphi^E_1}E_0\xrightarrow{\varphi^E_0}\Phi(E_1)\Bigr)
\end{gather*}
in $\mathcal{E}$ such that $\varphi^E_0\circ\varphi^E_1=w(E_1)$ and $\Phi\big(\varphi^E_1\big)\circ\varphi^E_0=w(E_0)$. Objects $E_1$ and $E_0$ in the above sequence are called the {\it components} of $E$. \end{dfn}


\begin{dfn}
Let $(\scrE,\Phi,w)$ be the above triple.
\begin{itemize}\itemsep=0pt
\item[$1.$] For two factorizations $E,F$ of $(\scrE,\Phi,w)$, we define the complex $\Hom(E,F)^{\bullet}$ with differential $d^{\bullet}_{(E,F)}\colon\Hom(E,F)^{\bullet}\to\Hom(E,F)^{\bullet+1}$ by
\begin{gather*}
{\rm Hom}(E,F)^{2n}\defeq{\rm Hom}_{\scrE}(E_1,\Phi^n(F_1))\oplus{\rm Hom}_{\scrE}(E_0,\Phi^n(F_0)),
\\
{\rm Hom}(E,F)^{2n+1}\defeq{\rm Hom}_{\scrE}(E_1,\Phi^{n}(F_0))\oplus{\rm Hom}_{\scrE}(E_0,\Phi^{n+1}(F_1)),
\end{gather*}
where $E_i$ and $F_i$ are the components of $E$ and $F$ respectively, and
\begin{gather*}
d^{\bullet}_{(E,F)}(f):=\varphi^F\circ f-(-1)^{{\rm deg}(f)}f\circ \varphi^E
\qquad\text{if}\quad f\in{\rm Hom}(E,F)^{{\rm deg}(f)}.
\end{gather*}
This defines the dg category
\begin{gather*}
\dgFact(\scrE,\Phi,w)
\end{gather*}
of factorizations of $(\scrE,\Phi,w)$.

\item[$2.$] The dg category $\dgFact(\scrE,\Phi,w)$ defines the additive categories
\begin{gather*}
\Fact(\scrE,\Phi,w)\defeq Z^0(\dgFact(\scrE,\Phi,w)),
\\
\K(\scrE,\Phi,w)\defeq H^0(\dgFact(\scrE,\Phi,w)).
\end{gather*}
We call $\K(\scrE,\Phi,w)$ the {\it homotopy category} of $(\scrE,\Phi,w)$.
\end{itemize}
\end{dfn}

The category $\Fact(\scrE,\Phi,w)$ is an exact category, and the homotopy category $\K(\scrE,\Phi,w)$ is a~triangulated category (see~\cite[Propositions~3.5 and~3.9]{H1}). To define the derived factorization categories, we define the totalizations of complexes of factorizations: For an object $E\in \Fact(\scrE,\Phi,w)$, let us set
\begin{gather*}
{\rm Com}(E)^{2i}\defeq\Phi^i(E_0), \qquad {\rm Com}(E)^{2i-1}\defeq\Phi^i(E_1),
\\
d_E^{2i}\defeq\Phi^i\big(\varphi_0^E\big),\qquad d_E^{2i-1}\defeq\Phi^i\big(\varphi_1^E\big).
\end{gather*}
Then a periodic (up to twists by $\Phi$) infinite sequence
\begin{gather*}
{\rm Com}(E)\defeq\big({\rm Com}(E)^{\bullet}, d_E^{\bullet}\big)
\end{gather*}
 in $\scrE$ satisfies $d_E^{i+1}\circ d_E^i=w({\rm Com}(E)^i)$ for all $i\in\bZ$.

\begin{dfn}
Let $E^{\bullet}=\big(\cdots\rightarrow E^i\xrightarrow{\delta^i}E^{i+1}\rightarrow\cdots\big)$ be a bounded complex of $\Fact(\scrE,\Phi,w)$. We define the {\it totalization} $\Tot(E^{\bullet})\in\Fact(\scrE,\Phi,w)$ of $E^{\bullet}$ by
\begin{gather*}
{\rm Tot}(E^{\bullet})\defeq\big(T_1\xrightarrow{t_1}T_0\xrightarrow{t_0}\Phi(T_1)\big),
\end{gather*}
where
\begin{gather*}
T_l:=\bigoplus_{i+j=-l}{\rm Com}(E^i)^j,
\\
t_l|_{{\rm Com}(E^i)^j}\defeq{\rm Com}(\delta^i)^j+(-1)^id_{E^i}^j.
\end{gather*}
Taking totalizations defines an exact functor
\begin{gather*}
\Tot\colon\ \Chb(\Fact(\scrE,\Phi,w))\rightarrow \Fact(\scrE,\Phi,w).
\end{gather*}
\end{dfn}

\begin{dfn}
Let $\A(\scrE,\Phi,w)$ be the smallest thick subcategory of $\K(\scrE,\Phi,w)$ that contains tota\-lizations of all short exact sequences in $\Fact(\scrE,\Phi,w)$. Then we define the {\it derived factorization category} $\D(\scrE,\Phi,w)$ of the triple $(\scrE,\Phi,w)$ by the Verdier quotient
\begin{gather*}
\D(\scrE,\Phi,w)\defeq\K(\scrE,\Phi,w)/\A(\scrE,\Phi,w).
\end{gather*}
If $\scrE$ is closed under coproducts in $\scrA$, we denote by
$\Aco(\scrE,\Phi,w)$ the smallest thick subcategory of $\K(\scrE,\Phi,w)$ that contains totalizations of all short exact sequences in $\Fact(\scrE,\Phi,w)$ and closed under coproducts. Then we define the {\it coderived factorization category} $\Dco(\scrE,\Phi,w)$ by
\begin{gather*}
\Dco(\scrE,\Phi,w)\defeq\K(\scrE,\Phi,w)/\Aco(\scrE,\Phi,w).
\end{gather*}
\end{dfn}

Let $\scrF$ be an exact subcategory of another abelian category $\scrB$ satisfying the same properties of $\scrA$, and let
\begin{gather*}
(\scrF,\Psi,v)
\end{gather*} be a triple as in \eqref{triple}.

\begin{dfn} An additive functor
\begin{gather*}
F\colon\ \scrE \to \scrF
\end{gather*}
is {\it factored} with respect to the potentials $(\Phi,w)$ and $(\Psi,v)$ if there is a functor isomorphism
\begin{gather*}
\alpha\colon\ F\circ \Phi \simto \Psi\circ F
\end{gather*}
such that for every object $E\in \scrE$, the following diagram commutes:
\[
\begin{tikzcd}
&F(E)\arrow[ld, "F(w(E))"']\arrow[rd, "v(F(E))"]&\\
F(\Phi(E))\arrow[rr, "\alpha(E)"]&& \Psi(F(E)).
\end{tikzcd}
\]
\end{dfn}

If $F\colon \scrE \to \scrF$ is a factored functor with respect to $(\Phi,w)$ and $(\Psi,v)$, then it induces a dg functor
\begin{gather*}
F\colon\ \dgFact(\scrE,\Phi,w)\to\dgFact(\scrF,\Psi,v).
\end{gather*}
This dg functor defines the additive functors
\begin{gather*}
F\colon\ \Fact(\scrE,\Phi,w)\to\Fact(\scrF,\Psi,v),
\\
F\colon\ \K(\scrE,\Phi,w)\to\K(\scrF,\Psi,v),
\end{gather*}
where the latter functor is an exact functor. To define the derived functors of exact functors between homotopy categories, we need the following results due to~\cite{bdfik}.

\begin{prop}[{\cite[Corollary~2.25]{bdfik}}]\label{inj resol}
Assume that $\scrA$ has enough injectives and that the coproducts of injectives are injective. Let $\scrI\subset \scrA$ be the subcategory of injective objects. Then the natural functor
\begin{gather*}
\K(\scrI,\Phi,w)\to\Dco(\scrA,\Phi,w)
\end{gather*}
is an equivalence.
\end{prop}

\begin{prop}\label{proj resol}
Let $\scrP\subset \scrA$ be the subcategory of projective objects in $\scrA$, and assume that~$\scrA$ has enough projectives. Let $\scrC\subset \scrA$ be an abelian subcategory that is preserved by $\Phi$, and let $\scrQ\defeq\scrC\cap \scrP$.
\begin{itemize}\itemsep=0pt
\item[$1.$] Assume that all objects in $\scrC$ are compact in $\scrA$. Then for any $P\in \K(\scrQ,\Phi,w)$ and $A\in \Aco(\scrA,\Phi,w)$ we have
\begin{gather*}
\Hom_{\K(\scrA,\Phi,w)}(P,A)=0.
\end{gather*}
In particular,
the natural functor
\begin{gather*}
\K(\scrQ,\Phi,w)\to\Dco(\scrA,\Phi,w)
\end{gather*}
is fully faithful.

\item[$2.$] If every object in $\scrC$ has a finite projective resolution in $\scrC$, the natural functor
\begin{gather*}
\K(\scrQ,\Phi,w)\to\D(\scrC,\Phi,w)
\end{gather*}
is essentially surjective.
\end{itemize}
\begin{proof}
(1) The latter statement follows from the former one by~\cite[Proposition~B.2]{ls}. The~vani\-shing $\Hom_{\K(\scrA,\Phi,w)}(P,A)=0$ reduces to the case when $A\in \A(\scrA,\Phi,w)$, since the components of~$P$ are compact by our assumption. If $A\in \A(\scrA,\Phi,w)$, the vanishing follows from~\cite[Lemma~2.4]{bdfik}.

(2) This follows from~\cite[Proposition~2.22]{bdfik}.
\end{proof}
\end{prop}

\begin{dfn} Let $F\colon \scrA\to \scrB$ be a factored functor with respect to $(\Phi,w)$ and $(\Psi,v)$.

\begin{itemize}\itemsep=0pt
\item[1.] Assume that $\scrA$ has enough injectives and that the coproducts of injectives are injective.
 If $F$ is left exact, we define the {\it right derived functor}
\begin{gather*}
\bR F\colon\ \Dco(\scrA,\Phi,w)\to \Dco(\scrB,\Psi,v)
\end{gather*}
of $F\colon \K(\scrA,\Phi, w)\to\K(\scrB,\Psi,v)$ to be the composition
\[
\begin{tikzcd}
\Dco(\scrA,\Phi,w)\arrow[r,"\sim"]&\K(\scrI,\Phi,w)\arrow[r,"F"]&\K(\scrB,\Psi,v)\arrow[r,"Q"]&\Dco(\scrB,\Psi,v),
\end{tikzcd}
\]
where the first functor is the equivalence in Proposition~\ref{inj resol}, and $Q$ is the natural quotient functor.

 \item[2.] Let $\scrC\subset \scrA$ and $\scrD\subset \scrB$ be abelian subcategories that are preserved by $\Phi$ and $\Psi$ respectively, and assume that the factored functor $F$ restricts to $F\colon \scrC\to \scrD$. Assume that every object in $\scrC$ is compact in $\scrA$ and has a finite projective resolution in $\scrC$, and that the natural functor $\D(\scrC,\Phi,w)\to \Dco(\scrA,\Phi,w)$ is fully faithful.
If $F$ is right exact, we define the {\it left derived functor}
\begin{gather*}
\bL F\colon\ \D(\scrC,\Phi,w)\to \D(\scrD,\Psi,v)
\end{gather*}
of $F\colon \K(\scrC,\Phi,w)\to \K(\scrD,\Psi,v)$ to be the composition
\[
\begin{tikzcd}
\D(\scrC,\Phi,w)\arrow[r,"\sim"]&\K(\scrQ,\Phi,w)\arrow[r,"F"]&\K(\scrD,\Psi,v)\arrow[r,"Q"]&\D(\scrD,\Psi,v),
\end{tikzcd}
\]
where the first functor is the equivalence in Proposition~\ref{proj resol}.
\end{itemize}
\end{dfn}

\subsection{Quasi-coherent modules over sheaves of algebras}

In this subsection, we recall fundamental properties of quasi-coherent modules over sheaves of~algebras over schemes.

Let $X$ be a scheme. An $\cO_X$-module $\cF$ is said to be {\it quasi-coherent} if for any point $x\in X$, there are an open neighborhood $U$ of $x$
and an exact sequence
\begin{gather*}
\cO_U^{\oplus I}\to \cO_U^{\oplus J}\to \cF|_U\to0,
\end{gather*}
where $I$ and $J$ are sets. A quasi-coherent $\cO_X$-module $\cF$ is called a {\it coherent} $\cO_X$-module, if for every affine open subset $U=\Spec R$ of $X$, $\Gamma(U,\cF)$ is a finitely generated $R$-module. We denote by $\Qcoh X$ (resp.\ $\coh X$) the category of quasi-coherent (resp.\ coherent) $\cO_X$-modules.

\begin{dfn}
An {\it $X$-algebra} is a sheaf $\cA$ of not necessarily commutative algebras together with a morphism $\uprho\colon \cO_X\to \cA$ of sheaves of algebras such that the image of $\uprho$ lies in the center of~$\cA$ and that $\cA$ is a quasi-coherent $\cO_X$-module. An $X$-algebra $\uprho\colon \cO_X\to \cA$ is said to be {\it coherent} if $\cA$ is a coherent $\cO_X$-module. A {\it morphism} of $X$-algebras is a morphism $\varphi\colon \cA\to \cB$ of~sheaves of algebras that is $\cO_X$-linear.

For an $X$-algebra $\uprho\colon \cO_X\to\cA$, a {\it quasi-coherent $($resp.\ coherent$)$ $\cA$-module} is a right $\cA$-mo\-dule~$\cM$ that is a quasi-coherent (resp.\ coherent) $\cO_X$-module. For quasi-coherent $\cA$-modu\-les~$\cM$ and~$\cN$, a {\it morphism} from $\cM$ to $\cN$ is a morphism of sheaves of right $\cA$-modules. We~denote~by\vspace{-.7ex}
\begin{gather*}
\Qcoh \cA
\end{gather*}
the category of quasi-coherent $\cA$-modules, and we write $\coh \cA$ for the full subcategory of coherent $\cA$-modules.
\end{dfn}

\begin{rem}%\label{ab5}
Since the category of right $\cA$-modules and the category $\Qcoh X$ are both abelian, $\Qcoh \cA$ is also an abelian category. Furthermore, in Proposition~\ref{grothen prop} we will see that $\Qcoh \cA$ is~a~Grothendieck category.
\end{rem}

Let $\varphi\colon \cA\to \cB$ be a morphism of $X$-algebras. For an object $\cN\in \Qcoh \cB$, we define a~qua\-si-coherent $\cA$-module $\cN_{\varphi}$ by the composition\vspace{-.7ex}
 \begin{gather*}
 \cN\times \cA\xrightarrow{\id\times \varphi} \cN\times \cB\to \cN,
 \end{gather*}
 where the second morphism is the right $\cB$-action of $\cN$. This defines the functor\vspace{-.7ex}
\begin{equation}\label{restriction}
(-)_{\varphi}\colon\ \Qcoh \cB\to \Qcoh \cA,
\end{equation}
and we call the functor $(-)_{\varphi}$ the {\it restriction by $\varphi$}.

Conversely, for an object $\cM\in \Qcoh \cA$, the tensor product
$\cM\otimes_{\cA}\cB$ has a natural right $\cB$-action, and this defines the functor\vspace{-.7ex}
\begin{equation}\label{extension}
(-)\otimes_\cA\cB\colon\ \Qcoh \cA\to \Qcoh \cB,
\end{equation}
and we call this functor the {\it extension by $\varphi$}.
It is standard that we have an adjunction\vspace{-.7ex}
\begin{gather*}
(-)\otimes_\cA\cB \dashv (-)_{\varphi}.
\end{gather*}

Although the following propositions might be well known, we give the proofs for reader's convenience.
\begin{prop}
Let $X$ be a scheme, and $\uprho\colon \cO_X\to\cA$ an $X$-algebra.\vspace{-.5ex}
\begin{itemize}\itemsep=-1pt
\item[$1.$] For every $\cM\in \Qcoh \cA$ and every point $x\in X$, there are an open neighborhood $U$ of $x$ and a surjective $\cA|_U$-linear map\vspace{-.7ex}
\begin{gather*}
\cA|_U^{\oplus J}\surjto \cM|_U,
\end{gather*}
where $J$ is a set.
\item[$2.$] A right $\cA$-module $\cM$ is quasi-coherent if and only if for every point $x\in X$ there are an~open neighborhood $U$ and an exact sequence\vspace{-.7ex}
 \begin{gather*}
 \cA|_U^{\oplus I}\to \cA|_U^{\oplus J}\to \cM|_U\to 0.
 \end{gather*}
 in the category of right $\cA$-modules.
\end{itemize}
 \begin{proof}
 (1) Let $\cM\in \Qcoh \cA$ and $x\in X$. Then there are an open neighborhood $U$ of $x$, a set $J$ and an surjective morphism
 $\pi\colon \cO_U^{\oplus J}\surjto\cM_{\uprho}|_U$
 in $\Qcoh X$. Since the functor $(-)\otimes_{\cO_U}\cA|_U\colon \Qcoh U\to \Qcoh \cA|_U$ has a right adjoint functor, it is right exact and commutes with small direct sums. Hence we have a natural isomorphism $\cO_U^{\oplus J}\otimes _{\cO_U}\cA_U\cong \cA|_U^{\oplus J}$ and the morphism
{\samepage \begin{gather*}
 \pi\otimes_{\cO_U}\cA|_U\colon\ \cA_U^{\oplus J}\to \cM_{\uprho}|_U\otimes_{\cO_U}\cA|_U
 \end{gather*}
 is surjective. Composing this surjection with a natural surjective morphism
 $\cM_{\uprho}|_U\otimes_{\cO_U}\cA|_U\surjto \cM|_U$; $m\otimes a\mapsto ma$, we have a surjective morphism $\cA|_U^{\oplus J}\surjto \cM|_U$.}

 (2) ($\Rightarrow$) Assume that a right $\cA$-module $\cM$ is quasi-coherent and let $x\in X$ be a point. Then by (1) there are an open neighbborhood $U$ of $x$ and a surjective $\cA|_U$-linear map $p \colon \cA|_U^{\oplus J}\surjto \cM|_U$. Denote by $\cK$ the Kernel of $p$ and by $i\colon \cK\hookto \cA|_U^{\oplus J}$ the natural inclusion. Then $\cK$ is also a quasi-coherent $\cA$-module, and thus, by shrinking $U$ if necessary, we have a surjective $\cA|_U$-linear map $q \colon \cA|_U^{\oplus I}\surjto \cK|_U$. Then the sequence
 \begin{gather*}
 \cA|_U^{\oplus I}\xrightarrow{i \circ q}\cA|_U^{\oplus J}\xrightarrow{p}\cM|_U\to 0
 \end{gather*}
is an exact.

($\Leftarrow$) Assume that a right $\cA$-module $\cM$ is locally isomorphic to the cokernel of a morphism between free modules. Then, since $\cA$ is a quasi-coherent $\cO_X$-module and $\cM$ is locally isomorphic to the cokernel of a morphism of quasi-coherent modules, $\cM$ is a locally quasi-coherent module. Thus $\cM$ is a quasi-coherent $\cO_X$-module.
 \end{proof}
\end{prop}

\subsection{Equivariant sheaves}\label{equiv section}
We briefly recall the basics of equivariant quasi-coherent sheaves. For more details, see, for example,~\cite[Section~2]{bfk}.

Let $G$ be an algebraic group, and denote by $\varepsilon\colon \Spec k\to G$, $\mu\colon G\times G\to G$ and $\iota\colon G\to G$ the morphisms defining the identity, the multiplication and the inversion of $G$.
Let $X$ be a scheme with an algebraic left $G$-action $\sigma\colon G\times X \to X$. We denote by $\pi\colon G\times X\to X$ the natural projection, and we write
\begin{gather*}
\varepsilon_X\colon\ X\to G\times X
\end{gather*}
for the composition $X\simto \Spec k\times X\xrightarrow{ \varepsilon\times \id_X}G\times X$. \begin{dfn}
A {\it $G$-equivariant structure} on a quasi-coherent sheaf $\cF\in \Qcoh X$ is an isomorphism
\begin{gather*}
\uptheta\colon\ \pi^*\cF\simto \sigma^*\cF
\end{gather*}
of $\cO_{G\times X}$-modules such that the equations
\begin{equation}\label{compati}
\varphi^*\uptheta\circ (\id_G\times\pi)^*\uptheta=(\mu\times\id_X )^*\uptheta \qquad\text{and}\qquad
\varepsilon_X^*\uptheta=\id_{\cF}
\end{equation}
hold, where $\varphi\colon G\times G\times X\to G\times X$ is the morphism defined by $\varphi(g,h,x)\defeq(h,gx)$.
\end{dfn}

\begin{rem}\label{closed point isom}
For any closed point $g\in G$, there is the corresponding morphism $g\colon \Spec k\to G$, and this induces the morphism $g\times\id_X\colon X\to G\times X$. The pull-back of a $G$-equivariant structure $\uptheta\colon \pi^*\cF\simto \sigma^*\cF$ by $g\times \id_X$ defines the isomorphism
\begin{gather*}
\uptheta_g\colon\ \cF\simto \sigma_g^*\cF,
\end{gather*}
where $\sigma_g\colon X\to X$ is the isomorphism defined by $\sigma_g\defeq \sigma\circ (g\times \id_X)$. The pull-back of the first equation in (\ref{compati}) by the morphism $ g\times h\times\id_X\colon X\to G\times G\times X$ implies the equation
\begin{gather*}
\sigma_g^*\uptheta_h\circ \uptheta_g=\uptheta_{gh},
\end{gather*}
and the pull-back of the second equation is nothing but the equation $\uptheta_{1_G}=\id_{\cF}$.
\end{rem}

A {\it $G$-equivariant quasi-coherent sheaf} \,is a pair $(\cF, \uptheta)$ of a quasi-coherent sheaf $\cF\in\Qcoh X$ and a $G$-equivariant structure $\uptheta$ on $\cF$. We say that $(\cF,\uptheta)$ is {\it coherent} (resp.\ {\it locally free}) if $\cF$ is coherent (resp.\ locally free). If there is no risk of confusion, a $G$-equivariant quasi-coherent sheaf $(\cF, \uptheta)$ is denoted simply by $\cF$.

Let $\big(\cF,\uptheta^{\cF}\big)$ and $\big(\cG,\uptheta^{\cG}\big)$ be $G$-equivariant quasi-coherent sheaves on $X$. A morphism $\varphi$: $\cF\to \cG$ of $\cO_X$-modules is said to be {\it $G$-equivariant} if the diagram
\begin{equation}\label{diag}
\begin{tikzcd}
\pi^*\cF\arrow[rr,"\pi^*\varphi"]\arrow[d,"\uptheta^{\cF}"']&&\pi^*\cG\arrow[d," \uptheta^{\cG}"]\\
\sigma^*\cF\arrow[rr, "\sigma^*\varphi"]&&\sigma^*\cG
\end{tikzcd}\end{equation}
is commutative.

\begin{rem}\label{cl pt}
Let $\big(\cF,\uptheta^{\cF}\big)$ and $\big(\cG,\uptheta^{\cG}\big)$ be $G$-equivariant quasi-coherent sheaves on $X$. Assume that $X$ is of finite type over $k$. Then the set of closed points of $G\times X$ is equal to the set of pairs $(x,g)$ of closed points $x\in X$ and $g\in G$. This implies that a morphism $\varphi\colon\cF\to \cG$ of the $\cO_X$-modules is $G$-equivariant if and only if for any closed point $g\in G$ the pull-back
\begin{equation*}
\begin{tikzcd}
\cF\arrow[rr,"\varphi"]\arrow[d,"\uptheta^{\cF}_g"']&&\cG\arrow[d," \uptheta^{\cG}_g"]\\
\sigma_g^*\cF\arrow[rr, "\sigma^*_g\varphi"]&&\sigma_g^*\cG
\end{tikzcd}
\end{equation*}
of the diagram (\ref{diag}) by $g\times\id_X $ is commutative.
\end{rem}

$G$-equivariant quasi-coherent sheaves on $X$ and $G$-equivariant morphisms define the abelian category
\begin{gather*}
\Qcoh_GX.
\end{gather*}
It is standard that the category $\Qcoh_GX$ is equivalent to the category $\Qcoh [X/G]$ of quasi-coherent sheaves on the quotient stack $[X/G]$;
\begin{equation}\label{equiv stack}
\Qcoh_GX\cong \Qcoh [X/G],
\end{equation}
 and for $\cF\in \Qcoh_GX$ we denote by $[\cF/G]\in \Qcoh [X/G]$ the image of $\cF$ by the equivalence \eqref{equiv stack}. We denote by $\coh_GX$ the subcategory of $G$-equivariant coherent sheaves on $X$. This category is an exact subcategory, and if $X$ is noetherian, it is an abelian subcategory. Furthermore, if $X$ is of finite type over $k$, by Remark~\ref{cl pt}, we have
\begin{equation}\label{G hom}
\Hom_{\Qcoh_GX}\big(\big(\cF,\uptheta^{\cF}\big),\big(\cG,\uptheta^{\cG}\big)\big) =\Hom_{\Qcoh X}(\cF,\cG)^G,
\end{equation}
where the right hand side $\Hom_{\Qcoh X}(\cF,\cG)^G$ is the $G$-invariant subset with respect to the $G$-action on $\Hom_{\Qcoh X}(\cF,\cG)$ defined by
$g\cdot\varphi\defeq (\uptheta^{\cG}_g)^{-1}\circ\sigma_g^*\varphi\circ\uptheta^{\cF}_g$.

Let $Y$ be another $G$-scheme, and $f\colon X\to Y$ a $G$-equivariant morphism that is quasi-compact and quasi-separated. Then $f$ induces the direct image\vspace{-.5ex}
\begin{gather*}
f_*\colon\ \Qcoh_G X\to \Qcoh_G Y
\end{gather*}
and the inverse image\vspace{-.5ex}
\begin{gather*}
f^*\colon\ \Qcoh_G Y\to \Qcoh_GX.
\end{gather*}
It is standard that $f^*$ is a left adjoint functor of $f_*$.

Tensor products and sheaf Hom define bi-functors
\begin{gather*}
(-)\otimes_{X}(-)\colon\ \Qcoh_G X\times \Qcoh_G X\to \Qcoh_G X,
\\[.5ex]
\cHom(-,-)\colon\ (\coh_G X)^{\rm op}\times \Qcoh_G X\to \Qcoh_G X,
\end{gather*}
and it is also standard that for any $\cF\in \coh_GX$, the functor $\cHom(\cF,-)\colon\Qcoh_G X\to \Qcoh_G X$ is right adjoint to the functor $(-)\otimes_{X}\cF\colon\Qcoh_G X\to \Qcoh_G X$.

\subsection{Tilting objects and derived equivalences}

We recall that tilting objects induce derived equivalences.
Throughout this subsection, $\scrA$ is an abelian category with small coproducts and enough injectives, and we assume that small coproducts of families of short exact sequences in $\scrA$ are exact. These properties of $\scrA$ are satisfied if $\scrA$ is a Grothendieck category.

\begin{dfn} Let $T\in \scrA$ be an object. The object $T$ is called a {\it tilting object} if the following conditions are satisfied:
\begin{itemize}\itemsep=0pt
\item[$1.$] $\Ext_{\scrA}^i(T,T)=0$ for all $i>0$.

\item[$2.$] $T$ is {\it compact} in $\D(\scrA)$, i.e., the natural map
\begin{gather*}
\bigoplus_{i\in \cI}\Hom_{\D(\scrA)}(T,A_i^{\bullet})\to\Hom_{\D(\scrA)}\Big(T,\bigoplus_{i\in\cI}A_i^{\bullet}\Big)
\end{gather*}
is an isomorphism of abelian groups for any set $\cI$ and any family $\{A_i^{\bullet}\}_{i\in I}$.

